%! Matrix Computations and A = LU — Unit 2, Lesson 2.3 — Warm-Up
\documentclass[10pt]{article}
\usepackage{linalg-article}
\usepackage{linalg-boxes}
\newcommand{\TallMath}[1]{$\displaystyle #1\rule[-1.4em]{0pt}{3.2em}$}
\newcommand{\xx}{\mathbf{x}}
\newcommand{\bb}{\mathbf{b}}
\newcommand{\cc}{\mathbf{c}}

\begin{document}

\pageheader{Unit 2, Lesson 2.3}{Warm-Up}
\namedateperiod

\vspace{0.12in}

\begin{notesbox}{Getting Ready: Skills We'll Use Today}
\small Today we \emph{save} elimination as a product $A=LU$ (lower times upper). These three moves are
the pieces. Answer quickly.

\vspace{4pt}
\textbf{1. Lower times upper} (Lesson 1.4). Multiply, one column at a time:
\[
  \begin{bmatrix} 1 & 0 \\ 4 & 1 \end{bmatrix}
  \begin{bmatrix} 2 & 5 \\ 0 & 3 \end{bmatrix}
  = \begin{bmatrix}\rule{0pt}{1.1em}\blank{1.1cm} & \blank{1.1cm}\\[2pt]\blank{1.1cm} & \blank{1.1cm}\end{bmatrix}.
\]
Notice a lower-triangular times an upper-triangular matrix rebuilds a \emph{full} matrix.

\vspace{4pt}
\textbf{2. Back-substitution} (Lesson 2.1). The system is already upper-triangular. Solve it from the
\emph{bottom row up}:
\[
  \begin{aligned}
    x + 2y &= 5 \\
    y &= 2
  \end{aligned}
  \qquad\Longrightarrow\qquad y = \blank{1.1cm}, \quad x = \blank{1.1cm}.
\]

\vspace{4pt}
\textbf{3. One multiplier} (Lesson 2.1). To clear the $x$ below the pivot in
$\ x + 3y = 7,\ \ 2x + 7y = 15$, the pivot is $1$ and the multiplier is
\[
  \ell = \frac{\text{entry to eliminate}}{\text{pivot}} = \frac{2}{1} = \blank{1.2cm}.
\]
\end{notesbox}

\end{document}
